% ======================================================================
% section_theory.tex — Budget-of-uncertainty robust workload constraint
% for the CSLAP placement model. THEORY ONLY (no implementation).
%
% Compilable inside a standard article-class document with
% amsmath, amssymb, amsthm, natbib (authoryear), and cleveref
% (\usepackage[capitalize,noabbrev,nameinlink]{cleveref}, as in the
% project drafts) loaded, and with
%   \newtheorem{proposition}{Proposition}
%   \newtheorem{remark}{Remark}
% defined in the preamble. Citations use \citet/\citep with the keys
% bertsimas2004 (the robust draft's existing key for Bertsimas & Sim 2004;
% do NOT introduce a duplicate bertsimas2004price entry on merge),
% bertsimas2003robust (Bertsimas & Sim 2003, Math. Program.),
% soyster1973convex, bental2009robust.
% The label sec:empirical is wired by the assembler.
% ======================================================================

\section{A budget-of-uncertainty robust counterpart for the workload constraint}
\label{sec:theory}

The empirical companion (\Cref{sec:empirical}) documents a failure mode of
trained CSLAP layouts under demand drift: layouts that satisfy the station
workload constraint on the training window violate it out-of-sample on early
temporal cuts. This section develops the exact robust counterpart of the
workload constraint under budgeted pick-line uncertainty, targeting this
failure mode; how much protection is actually delivered depends on the
calibrated uncertainty set and is assessed in \Cref{sec:calibration}.
Uncertainty is placed on the pick-line coefficients $L_p$ only; the visits
objective is untouched. The construction is the budget-of-uncertainty model
of \citet{bertsimas2004}, in the row-wise robust-optimization framework of
\citet{bental2009robust}, specialized to the CSLAP placement MILP.
Symbols follow the project notation: $P$, $O$, $S$ are the product, order, and
station sets; $P_o$ the support of order $o$; $\zeta_s$, $V_s$, $T_s$ the slot
capacity, processing speed, and time capacity of station $s$.

\subsection{Nominal placement model}
\label{sec:nominal}

% --- removable block: restates the nominal model of the COR manuscript with
% --- its canonical labels (obj, con14--con17) so this section compiles
% --- standalone. Delete on merge into the manuscript and keep the \eqref's.
The nominal placement model assigns every product to exactly one station and
minimizes the total number of station visits over the training orders. The
decision variables are $x_{ps} \in \{0,1\}$ (product $p \in P$ assigned to
station $s \in S$) and the auxiliary visit variables $z_{os} \in [0,1]$ (order
$o \in O$ visits station $s$, completely dependent on $x_{ps}$):
\begin{equation}
\min \sum_{o \in O} \sum_{s \in S} z_{os}
\label{obj}
\end{equation}
subject to
\begin{align}
&\sum_{s \in S} x_{ps} = 1 \quad \forall p \in P, \label{con14}\\
&\sum_{p \in P} x_{ps} \leq \zeta_s \quad \forall s \in S, \label{con15}\\
&x_{ps} \leq z_{os} \quad \forall p \in P_o,\ \forall o \in O,\ \forall s \in S, \label{con16}\\
&\sum_{p \in P} x_{ps} \cdot \frac{L_p}{V_s} \leq T_s \quad \forall s \in S, \label{con17}\\
&x_{ps} \in \{0, 1\}, \quad z_{os} \in [0,1], \quad \forall p \in P,\ o \in O,\ s \in S. \nonumber
\end{align}
Constraint \eqref{con14} assigns each product to exactly one station,
\eqref{con15} is the physical slot capacity, \eqref{con16} links assignments to
visits (a logical OR; no big-$M$ is needed because the objective pushes
$z_{os}$ down to $\max_{p \in P_o} x_{ps}$), and \eqref{con17} bounds the
workload of each station over the planning window.\footnote{The manuscripts typeset the summation in
\eqref{con17} as $\sum_{p \in P_o}$ with no quantifier on $o$; the column-generation
pricing model and the implementation both sum over \emph{all} products assigned
to the station, so we write $\sum_{p \in P}$ here.}

The workload coefficient $L_p$ is a \emph{point estimate}: the training-window
pick-line count of product $p$, computed as the number of distinct training
orders containing $p$ (\texttt{REAL\_LINES}; on instances with duplicate
order--product rows the raw-line and distinct-order conventions differ ---
see the BERNER caveat in \Cref{sec:empirical}). The time capacity is
calibrated on the same window by the
project's $10\%$ operational-slack rule,
\begin{equation}
T_s \;=\; \Bigl\lceil\, 1.10 \cdot \frac{\sum_{p \in P} L_p}{V_s\,|S|} \,\Bigr\rceil,
\label{eq:bs-slack}
\end{equation}
so a layout solving \eqref{obj}--\eqref{con17} is guaranteed workload-feasible
\emph{only for the estimate $L_p$ of the training window}. When the realized
pick-line vector of a later operating window drifts away from the estimate ---
as the temporal experiments of \Cref{sec:empirical} show it does ---
\eqref{con17} offers no protection: a single product whose demand doubles can
push its station past $T_s$ while the model still certifies feasibility.

\subsection{Budget-of-uncertainty model for pick-line workloads}
\label{sec:budget}

Let $L^0_p$ denote the nominal (training-window) estimate of the pick-line
count --- the quantity written $L_p$ in \eqref{con17} --- and let
$\hat L_p \ge 0$ be a deviation magnitude for product
$p$.\footnote{The superscripted symbol $L^0_p$ is chosen deliberately: the
accented $\bar L_p$ is reserved by the covering draft for the
closure-inflated line count (and $\bar T_s$ for its recalibrated capacity),
so no renaming rule is needed when the two studies are merged into one
manuscript.} The uncertain coefficient is modeled as
\begin{equation}
\tilde L_p \in \bigl[L^0_p - \hat L_p,\ L^0_p + \hat L_p\bigr]
\quad \forall p \in P,
\label{eq:bs-uncset}
\end{equation}
with the budget restriction of \citet{bertsimas2004}: in any realization,
\emph{at most $\Gamma$
products deviate from their nominal value simultaneously}, where
$\Gamma \in \{0, 1, \dots, |P|\}$ is the conservatism knob. Since the workload
constraint is a $\le$ row with nonnegative coefficients, only upward
deviations threaten it, so the worst case places every deviating coefficient
at its upper endpoint $L^0_p + \hat L_p$.

Write $\Phi_s = \{p \in P : x_{ps} = 1\}$ for the set of products assigned to
station $s$ (the Hexaly set-variable of the project notation). Requiring
\eqref{con17} to hold for every realization in the budgeted set gives the
robust workload constraint
\begin{equation}
\sum_{p \in \Phi_s} \frac{L^0_p}{V_s}\, x_{ps}
\;+\;
\max_{\substack{J \subseteq \Phi_s \\ |J| \le \Gamma}}\;
\sum_{p \in J} \frac{\hat L_p}{V_s}\, x_{ps}
\;\le\; T_s
\qquad \forall s \in S.
\label{eq:bs-robust}
\end{equation}
On $\Phi_s$ the factors $x_{ps} = 1$ are redundant; they are retained to make
the decision-dependence of the adversary's ground set explicit, and they let
the sums range over all of $P$ without changing value (terms with
$x_{ps} = 0$ vanish). The inner combinatorial maximization admits the
equivalent continuous relaxation
\begin{equation}
\psi_s(x;\Gamma) \;=\; \max_{z} \;
\sum_{p \in P} \frac{\hat L_p}{V_s}\, x_{ps}\, z_p
\quad \text{s.t.} \quad
\sum_{p \in P} z_p \le \Gamma, \qquad
0 \le z_p \le 1 \ \ \forall p \in P.
\label{eq:bs-inner}
\end{equation}
For fixed $x$, \eqref{eq:bs-inner} is a linear program in $z$ alone: the $x$
enter only through the constant objective coefficients
$d_{ps} := (\hat L_p / V_s)\, x_{ps} \ge 0$. Its constraint matrix --- one
cardinality row plus the identity box --- is an interval matrix, hence totally
unimodular; equivalently, by the greedy argument, since all $d_{ps} \ge 0$ an
optimal $z$ sets $z_p = 1$ on the $\Gamma$ largest coefficients (on all
positive ones if fewer than $\Gamma$) and $z_p = 0$ elsewhere. For integer
$\Gamma$ integrality of the inner adversarial problem is therefore
\emph{free}, and $\psi_s(x;\Gamma)$ equals the set-maximum in
\eqref{eq:bs-robust}.\footnote{The adversary variable $z_p$ (single subscript)
is unrelated to the visit variable $z_{os}$ and to the cover-pricing $z_p$ of
the covering draft; it exists only inside \eqref{eq:bs-inner} and is dualized
away in \Cref{prop:bs-counterpart}.}

\begin{proposition}[Exact linear counterpart]
\label{prop:bs-counterpart}
Fix a station $s \in S$, an integer budget $\Gamma \in \mathbb Z_{\ge 0}$,
and a binary assignment vector
$x_{\cdot s} \in \{0,1\}^{P}$.\footnote{Integrality of $\Gamma$ is what makes
the LP value $\psi_s(x;\Gamma)$ of \eqref{eq:bs-inner} coincide with the
set-maximum in \eqref{eq:bs-robust}; for fractional $\Gamma$ the
($\Rightarrow$) direction fails: one assigned product with
$\hat L_p / V_s = 1$, zero nominal slack, and $\Gamma = 1/2$ satisfies
\eqref{eq:bs-robust} (only $J = \emptyset$ is admissible), yet
$\psi_s(x;1/2) = 1/2 > 0$, so no multipliers exist. Real-valued budgets can
be recovered by adding the Bertsimas--Sim fractional deviation term
($\lfloor \Gamma \rfloor$ full deviations plus one scaled by
$\Gamma - \lfloor \Gamma \rfloor$) to \eqref{eq:bs-robust}; integer budgets
suffice here, consistent with $\Gamma \in \{0, 1, \dots, |P|\}$ above.}
The robust workload constraint
\eqref{eq:bs-robust} holds if and only if there exist multipliers
$\theta_s \ge 0$ and $\mu_{ps} \ge 0$, $p \in P$, such that
\begin{align}
&\sum_{p \in P} \frac{L^0_p}{V_s}\, x_{ps}
\;+\; \Gamma\, \theta_s
\;+\; \sum_{p \in P} \mu_{ps}
\;\le\; T_s, \label{eq:bs-cpt-load}\\[2pt]
&\theta_s + \mu_{ps} \;\ge\; \frac{\hat L_p}{V_s}\, x_{ps}
\qquad \forall p \in P. \label{eq:bs-cpt-dualfeas}
\end{align}
Since $\hat L_p / V_s$ and $L^0_p / V_s$ are data constants, each right-hand
side is a constant multiple of the decision variable $x_{ps}$, so
\eqref{eq:bs-cpt-load} and \eqref{eq:bs-cpt-dualfeas} are linear in
$(x, \theta, \mu)$; no integrality of $x$ is needed for linearity.
\end{proposition}

\begin{proof}
By the discussion following \eqref{eq:bs-inner} --- the step in which
integrality of $\Gamma$ is used --- constraint
\eqref{eq:bs-robust} is equivalent to
\begin{equation}
\psi_s(x;\Gamma) \;\le\; T_s - \sum_{p \in P} \frac{L^0_p}{V_s}\, x_{ps}
\;=:\; \rho_s(x),
\label{eq:bs-slackform}
\end{equation}
where $\rho_s(x)$ is the nominal slack of station $s$. The inner LP
\eqref{eq:bs-inner}, with coefficients
$d_{ps} = (\hat L_p / V_s)\, x_{ps} \ge 0$, is \emph{feasible} ($z = 0$) and
\emph{bounded above} ($\sum_p d_{ps} z_p \le \sum_p d_{ps} < \infty$ since
$z_p \le 1$), so it attains a finite optimum $\psi_s(x;\Gamma)$. Associate
the dual variable $\theta_s \ge 0$ with the budget row
$\sum_p z_p \le \Gamma$ and $\mu_{ps} \ge 0$ with each upper bound
$z_p \le 1$. The LP dual of \eqref{eq:bs-inner} is
\begin{equation}
\min_{\theta_s,\, \mu_{\cdot s}} \;\; \Gamma\,\theta_s + \sum_{p \in P} \mu_{ps}
\quad \text{s.t.} \quad
\theta_s + \mu_{ps} \ge d_{ps} \ \ \forall p \in P, \qquad
\theta_s \ge 0,\ \mu_{ps} \ge 0.
\label{eq:bs-dual}
\end{equation}
Because the primal is feasible and bounded, LP strong duality applies: the
dual \eqref{eq:bs-dual} is solvable and its optimal value equals
$\psi_s(x;\Gamma)$. We prove the two implications separately.

($\Leftarrow$) Suppose $(\theta_s, \mu_{\cdot s}) \ge 0$ satisfies
\eqref{eq:bs-cpt-load}--\eqref{eq:bs-cpt-dualfeas}. Constraint
\eqref{eq:bs-cpt-dualfeas} says the point is feasible for the dual
\eqref{eq:bs-dual}, so by \emph{weak} duality
$\psi_s(x;\Gamma) \le \Gamma\,\theta_s + \sum_p \mu_{ps}$. Rearranging
\eqref{eq:bs-cpt-load} gives
$\Gamma\,\theta_s + \sum_p \mu_{ps} \le \rho_s(x)$. Chaining the two
inequalities yields $\psi_s(x;\Gamma) \le \rho_s(x)$, which is
\eqref{eq:bs-slackform}, i.e., the robust constraint \eqref{eq:bs-robust}
holds.

($\Rightarrow$) Conversely, suppose \eqref{eq:bs-robust} holds, i.e.,
$\psi_s(x;\Gamma) \le \rho_s(x)$. By strong duality the dual
\eqref{eq:bs-dual} attains its optimum at some
$(\theta_s^\ast, \mu_{\cdot s}^\ast) \ge 0$ with
\[
\Gamma\,\theta_s^\ast + \sum_{p \in P} \mu_{ps}^\ast
\;=\; \psi_s(x;\Gamma) \;\le\; \rho_s(x)
\;=\; T_s - \sum_{p \in P} \frac{L^0_p}{V_s}\, x_{ps}.
\]
The optimal dual point is dual-feasible, so it satisfies
\eqref{eq:bs-cpt-dualfeas}; the displayed inequality is exactly
\eqref{eq:bs-cpt-load}. Hence the required multipliers exist.

Finally, linearity: $\hat L_p/V_s$ and $L^0_p/V_s$ are constants, so
\eqref{eq:bs-cpt-load} is linear in $(x,\theta,\mu)$ and the right-hand side
of \eqref{eq:bs-cpt-dualfeas} is a constant multiple of the decision variable
$x_{ps}$ --- no product of decision variables appears anywhere, whether or
not $x$ is integral.
\end{proof}

\subsection{The robust placement MILP}
\label{sec:robust-milp}

Applying \Cref{prop:bs-counterpart} station by station --- with a per-station
integer budget $\Gamma_s \in \mathbb Z_{\ge 0}$, the uniform choice
$\Gamma_s = \Gamma$ being the default --- and substituting the counterpart
for \eqref{con17} yields the full
robust placement MILP:
\begin{subequations}
\label{eq:bs-milp}
\begin{align}
&\min_{x,\,z,\,\theta,\,\mu} \; \sum_{o \in O} \sum_{s \in S} z_{os}
  \label{eq:bs-milp-obj}\\
\text{s.t.}\;\;
&\sum_{s \in S} x_{ps} = 1
  \quad \forall p \in P, \label{eq:bs-milp-assign}\\
&\sum_{p \in P} x_{ps} \le \zeta_s
  \quad \forall s \in S, \label{eq:bs-milp-slot}\\
&x_{ps} \le z_{os}
  \quad \forall p \in P_o,\ \forall o \in O,\ \forall s \in S, \label{eq:bs-milp-link}\\
&\sum_{p \in P} \frac{L^0_p}{V_s}\, x_{ps}
  + \Gamma_s\, \theta_s
  + \sum_{p \in P} \mu_{ps} \le T_s
  \quad \forall s \in S, \label{eq:bs-milp-wl}\\
&\theta_s + \mu_{ps} \ge \frac{\hat L_p}{V_s}\, x_{ps}
  \quad \forall p \in P,\ \forall s \in S, \label{eq:bs-milp-dualfeas}\\
&x_{ps} \in \{0,1\},\quad z_{os} \in [0,1],\quad
  \theta_s \ge 0,\quad \mu_{ps} \ge 0
  \quad \forall p,\ o,\ s. \label{eq:bs-milp-dom}
\end{align}
\end{subequations}
The objective \eqref{eq:bs-milp-obj} and the combinatorial core
\eqref{eq:bs-milp-assign}--\eqref{eq:bs-milp-link} are identical to the
nominal model \eqref{obj}, \eqref{con14}--\eqref{con16}; only the workload
block changes. The size overhead over the nominal MILP is exactly $+|S|$
continuous variables $\theta_s$, $+|P|\,|S|$ continuous variables $\mu_{ps}$,
and $+|P|\,|S|$ dual-feasibility constraints \eqref{eq:bs-milp-dualfeas};
the workload rows stay at $|S|$, the binaries and the objective are
unchanged. Robustness is bought with one continuous LP block, not with
additional integer structure.

\subsection{Properties}
\label{sec:properties}

\begin{proposition}[Budget limits and monotonicity]
\label{prop:bs-limits}
For an integer budget $\Gamma \in \mathbb Z_{\ge 0}$, let
\[
\mathcal A(\Gamma) \;=\; \bigl\{ x \in \{0,1\}^{P \times S} \;:\;
\exists\, (z, \theta, \mu) \text{ such that } (x, z, \theta, \mu)
\text{ is feasible for \eqref{eq:bs-milp} with } \Gamma_s = \Gamma \bigr\}
\]
denote the projection onto the assignment variables of the feasible set of
\eqref{eq:bs-milp} with uniform budget $\Gamma$. Then:
\begin{enumerate}
\item[(i)] $\Gamma = 0$ recovers the nominal model: $\mathcal A(0)$ equals
the projection onto $x$ of the feasible set of \eqref{con14}--\eqref{con17}
with $L_p = L^0_p$.
\item[(ii)] If $\Gamma \ge \zeta_s$ for every $s$, then \eqref{eq:bs-robust}
reduces to the box (Soyster) counterpart
$\sum_{p \in P} \frac{L^0_p + \hat L_p}{V_s}\, x_{ps} \le T_s$
\citep{soyster1973convex}.
\item[(iii)] $\Gamma \mapsto \mathcal A(\Gamma)$ is nonincreasing:
$\Gamma' \ge \Gamma$ implies $\mathcal A(\Gamma') \subseteq \mathcal A(\Gamma)$,
and the optimal visit count of \eqref{eq:bs-milp} is nondecreasing in $\Gamma$.
\end{enumerate}
\end{proposition}

\begin{proof}
(i) Both projections are cut out by the common core
\eqref{con14}--\eqref{con16}; the visit variables never restrict $x$ further
(setting $z_{os} = 1$ is always feasible), so it suffices to compare the
workload blocks. With $\Gamma = 0$ the inner maximum in \eqref{eq:bs-robust}
is over $J = \emptyset$ only and equals $0$, so \eqref{eq:bs-robust} is
\eqref{con17}. In the counterpart, given $x$ satisfying \eqref{con17}, choose
$\theta_s = \max_p \hat L_p / V_s$ and $\mu_{ps} = 0$:
\eqref{eq:bs-cpt-dualfeas} holds and $\Gamma_s \theta_s$ vanishes, so
\eqref{eq:bs-cpt-load} is exactly \eqref{con17}; conversely
\eqref{eq:bs-cpt-load} with $\Gamma_s = 0$ and $\mu \ge 0$ implies
\eqref{con17}.

(ii) By the slot capacity \eqref{con15}, $|\Phi_s| = \sum_p x_{ps} \le
\zeta_s \le \Gamma$, so the restriction $|J| \le \Gamma$ is vacuous and the
adversary deviates \emph{every} assigned product upward: the inner maximum is
$\sum_{p \in \Phi_s} \hat L_p / V_s$, the worst case of the full interval box
\eqref{eq:bs-uncset}. Any $\Gamma \ge \max_s |\Phi_s|$ suffices; since
$|\Phi_s|$ is decision-dependent, $\Gamma \ge \zeta_s$ is the natural
data-level sufficient condition.

(iii) $\psi_s(x;\Gamma)$ is nondecreasing in $\Gamma$ (the feasible set of
\eqref{eq:bs-inner} grows with $\Gamma$), so the feasible sets
$\mathcal A(\Gamma)$ are nested and minimizing the same objective over a
smaller set cannot decrease the optimum.
\end{proof}

Thus $\Gamma$ interpolates monotonically between the nominal model and full
box conservatism, exactly the Bertsimas--Sim design.

\begin{remark}[Decision-dependent adversary support]
\label{rem:bs-endogenous}
The adversary's ground set $\Phi_s$ in \eqref{eq:bs-robust} is itself a
decision, since $\Phi_s = \{p : x_{ps} = 1\}$. Endogenous uncertainty sets
are in general a hard extension of robust optimization, but no extra
machinery is needed here: because the factor $x_{ps}$ is kept inside the
inner objective of \eqref{eq:bs-inner}, deviating a non-assigned product
contributes zero, so the adversary restricted to $\Phi_s$ and the adversary
ranging over all of $P$ attain the same value. The dualization of
\Cref{prop:bs-counterpart} then applies verbatim. Linearity of the
counterpart is automatic because $\hat L_p / V_s$ are data constants (no
bilinear terms, no big-$M$); what binariness of $x_{ps}$ buys is this
ground-set equivalence, not linearity. (Under merely nonnegative $x_{ps}$
the equivalence holds with respect to the support $\{p : x_{ps} > 0\}$
rather than $\Phi_s$: a fractional $0 < x_{ps} < 1$ lies outside $\Phi_s$
yet contributes $(\hat L_p / V_s)\, x_{ps} > 0$ to \eqref{eq:bs-inner}.)
\end{remark}

\begin{remark}[Probabilistic guarantee]
\label{rem:bs-prob}
Suppose the realized deviations of the assigned products are independent and
symmetrically distributed on the intervals \eqref{eq:bs-uncset}. Then, for
any assignment $x$ satisfying the protected row \eqref{eq:bs-robust} with
budget $\Gamma$ --- in particular, any $x$ feasible for \eqref{eq:bs-milp}
--- \citet{bertsimas2004} bound the violation probability of that row by
\begin{equation}
\Pr\Bigl[\textstyle\sum_{p \in \Phi_s} \tilde L_p / V_s > T_s\Bigr]
\;\le\; B(n, \Gamma) \;\le\; \exp\!\bigl(-\Gamma^2 / (2n)\bigr),
\qquad n = |\Phi_s| \le \zeta_s,
\label{eq:bs-probbound}
\end{equation}
where $B(n, \Gamma)$ is the exact binomial bound of Theorem~2 of
\citet{bertsimas2004}, well approximated by
$1 - \Phi_{\mathcal N}\bigl((\Gamma - 1)/\sqrt{n}\bigr)$ with
$\Phi_{\mathcal N}$ the standard normal c.d.f.\ (not the assigned set
$\Phi_s$); the closed-form exponential expression is the weaker
Hoeffding-type bound of \citet{bertsimas2003robust}, derived for homogeneous
deviation magnitudes across the row. A budget of order $\sqrt{n}$ already
yields a strong guarantee at a small fraction of the box protection level.
One caveat: daily pick-line counts are nonnegative and right-skewed, so the
symmetric-deviation assumption behind \eqref{eq:bs-probbound} is an
approximation here; the bound is indicative, and in practice it is replaced
by the data-driven procedure of \Cref{sec:calibration}: deviation magnitudes
fitted on train-internal blocks, a budget fixed ex ante on a grid, and
retrospective diagnostics assessing adequacy.
\end{remark}

\begin{remark}[Constraint-side uncertainty only; relation to H1]
\label{rem:bs-objective}
Uncertainty enters model \eqref{eq:bs-milp} through the \emph{constraints}
only; the visits objective \eqref{eq:bs-milp-obj} is identical to the nominal
one. This is deliberately complementary to the falsified covering hypothesis
H1, which perturbed the \emph{objective} by enlarging every order to its
pattern-closure and paid a certified out-of-sample price ($G(\kappa) < 0$ at
$\kappa = 2$ on the clean instance). Here nothing is insured about the visit
count: the insured quantity is \emph{workload-feasibility risk} --- the
out-of-sample violations $W^{\mathrm n}_s > T_s$ documented in
\Cref{sec:empirical} ---
which the nominal model cannot see by construction. The two mechanisms are
orthogonal; the present study evaluates the workload channel alone.
\end{remark}

\begin{remark}[Volume basis and reporting convention]
\label{rem:bs-basis}
The counterpart is meaningful only if $L^0_p$, $\hat L_p$, and $T_s$ are
expressed on the \emph{same volume basis}. In the project's adapters, $T_s$
derives from the training-window totals by the $10\%$ slack rule
\eqref{eq:bs-slack} ($\kappa$-invariant on the clean instance); accordingly,
$L^0_p$ must be the training-window line count and $\hat L_p$ a deviation
of volume-normalized counts on that basis (\Cref{sec:calibration}). Mixing
bases --- e.g., a test-window $\hat L_p$ against a train-calibrated $T_s$ ---
would re-introduce the recalibration confound the clean instance was built to
remove. Since the protection level depends on $n = |\Phi_s| \le \zeta_s$, we
recommend reporting budgets as the dimensionless fraction
$\Gamma_s / \zeta_s$, comparable across stations and instances.
\end{remark}

\subsection{Data-driven calibration of $(\hat L_p, \Gamma)$}
\label{sec:calibration}

The counterpart of \Cref{sec:robust-milp} is fully specified once the
deviation magnitudes $\hat L_p$ and the budget $\Gamma$ are chosen. We define
the estimators the empirical section instantiates; no implementation is
discussed.

\paragraph{Deviation models.}
Let $\delta_p^{(b)} = L_p^{(b)} - L^0_p$ denote the realized deviation of
the volume-normalized line count $L_p^{(b)}$ of product $p$ on a
train-internal temporal block $b \in \{1,\dots,B\}$ (contiguous segments of
the training window by order rank; counts are rescaled so each block carries
the training-window total volume, keeping everything on the basis of
\Cref{rem:bs-basis}). The three deviation models are
\begin{equation}
\text{M1:}\ \ \hat L_p = c_1\, L^0_p,
\qquad
\text{M2:}\ \ \hat L_p = c_2\, \sqrt{L^0_p},
\qquad
\text{M3:}\ \ \hat L_p = \operatorname{std}_{b=1,\dots,B}\bigl(L_p^{(b)}\bigr).
\label{eq:bs-devmodels}
\end{equation}
M1 assumes deviations proportional to demand (multiplicative drift), M2
assumes the $\sqrt{L^0_p}$ scaling of count noise (Poisson-like), and M3
is the nonparametric per-product standard deviation across the training
blocks. M3 needs no fitted constant; M1 and M2 are calibrated as follows.

\paragraph{Fitting $c$.}
For a coverage target $q \in (0,1)$ (e.g.\ $q = 0.9$), the constant is the
smallest value whose implied intervals cover a $L^0$-weighted fraction $q$
of the realized deviations on the train-internal blocks:
\begin{equation}
c(q) \;=\; \min \Bigl\{ c \ge 0 \;:\;
\sum_{p,\,b} L^0_p\, \mathbf 1\bigl[\,|\delta_p^{(b)}| \le c\, g(L^0_p)\bigr]
\;\ge\; q \sum_{p,\,b} L^0_p \Bigr\},
\label{eq:bs-cfit}
\end{equation}
with $g(L^0_p) = L^0_p$ for M1 and $g(L^0_p) = \sqrt{L^0_p}$ for
M2; equivalently, $c(q)$ is the $L^0$-weighted $q$-quantile of the
normalized deviations $|\delta_p^{(b)}| / g(L^0_p)$. Weighting by
$L^0_p$ prioritizes high-volume products, whose deviations dominate
station workloads; restricting to train-internal blocks keeps the calibration
free of test leakage.

\paragraph{Budget diagnostics.}
Fix a trained layout $a$ (indicator $x^a$), a station $s$ with assigned set
$\Phi_s$, nominal load $\bar W_s = \sum_{p \in \Phi_s} L^0_p / V_s$, and
realized normalized load
$W^{\mathrm n}_s = \sum_{p \in \Phi_s} L^{\mathrm n}_p / V_s$, where
$L^{\mathrm n}_p$ is the volume-normalized realized line count of the
evaluation window. By \eqref{eq:bs-inner},
$\psi_s(x^a;\Gamma) = \frac{1}{V_s} \sum_{p \in J_\Gamma(s)} \hat L_p$ with
$J_\Gamma(s) \subseteq \Phi_s$ the $\min(\Gamma, |\Phi_s|)$ largest
deviations at station $s$. Two \emph{retrospective} diagnostics summarize
how large a budget the data ask for and how large a budget the training
instance tolerates:
\begin{align}
\Gamma^{\mathrm{cover}}(s) &\;=\;
\min \Bigl\{ \Gamma \in \mathbb Z_{\ge 0} \;:\;
\bar W_s + \psi_s(x^a;\Gamma) \;\ge\; W^{\mathrm n}_s \Bigr\},
\label{eq:bs-gcover}\\
\Gamma^{\mathrm{bind}}(s) &\;=\;
\min \Bigl\{ \Gamma \in \mathbb Z_{\ge 0} \;:\;
\bar W_s + \psi_s(x^a;\Gamma) \;>\; T_s \Bigr\}.
\label{eq:bs-gbind}
\end{align}
$\Gamma^{\mathrm{cover}}(s)$ is the smallest budget whose protection level
reaches the realized load of station $s$: the budget that would have been
needed for the robust constraint to ``see'' the observed drift (it is $0$
when $\bar W_s \ge W^{\mathrm n}_s$, and undefined --- reported as $\infty$
--- when even $\Gamma = |\Phi_s|$ under-covers, i.e., when the deviation
model itself is too small). $\Gamma^{\mathrm{bind}}(s)$ is the smallest
budget at which the robust left-hand side exceeds $T_s$ on the training data,
i.e., at which the counterpart starts to cut off the trained layout; budgets
below it exert no retraining pressure at station $s$, budgets at or above it
force the optimizer to rebalance. Symmetrically to
$\Gamma^{\mathrm{cover}}$, the set in \eqref{eq:bs-gbind} may be empty: when
even the full box load keeps the station within capacity,
$\bar W_s + \sum_{p \in \Phi_s} \hat L_p / V_s \le T_s$, no finite budget
makes the constraint bind ($\psi_s$ saturates at $\Gamma = |\Phi_s|$), and
we set $\Gamma^{\mathrm{bind}}(s) = \infty$ --- such stations tolerate any
budget.

\paragraph{Use in the empirical study.}
In a future implementation phase the operative budget would be fixed
\emph{ex ante}: robust layouts would be trained on a predeclared grid of
budgets (reported as fractions of $\zeta_s$ per \Cref{rem:bs-basis}), with
the deviation magnitudes fitted on train-internal blocks only via
\eqref{eq:bs-cfit}, so that no evaluation-window quantity would enter the
choice of $\Gamma$ for any scored test cut. No robust layout is trained in
this study: \Cref{sec:empirical} computes only the \emph{retrospective}
diagnostics \eqref{eq:bs-gcover}--\eqref{eq:bs-gbind} on the existing
nominal layouts --- for each temporal cut and deviation model in
\eqref{eq:bs-devmodels}, the distributions of $\Gamma^{\mathrm{cover}}(s)$
over the stations that actually violate $W^{\mathrm n}_s > T_s$
out-of-sample and of $\Gamma^{\mathrm{bind}}(s)$ over \emph{all} stations
(healthy and violated alike) --- answering ``what budget \emph{would have
been} needed'' and ``what budget the training instance tolerates''. The
empirical instantiation also departs from \eqref{eq:bs-cfit} on the
magnitude side: the reported constants $c_1, c_2$ (\Cref{sec:gamma}) apply
the same $L^0$-weighted quantile formula \emph{retrospectively} to the
realized evaluation-window deviations
$\delta_p = L^{\mathrm n}_p - L^0_p$ --- they are diagnostic, not
planner-available --- and only M3 uses train-internal blocks as
\eqref{eq:bs-cfit} prescribes. A budget $\Gamma$ would be judged
\emph{adequate in retrospect} when it dominates the observed
$\Gamma^{\mathrm{cover}}$ quantiles on drift-hit stations --- so the
protected constraint would have flagged the realized violations, and the
robust model is designed to rebalance them away at train time --- while
staying below $\Gamma^{\mathrm{bind}}$ on most healthy stations, so
protection is bought where drift occurs rather than by uniform slack. Once
an ex-ante grid exists, these curves would audit it after the fact rather
than select the budget: they replace the a priori appeal to the
symmetric-noise bound \eqref{eq:bs-probbound} as the lens for \emph{judging}
$\Gamma$, and they must not feed back into its choice on any window used
for scoring.
